% SOURCE_STATEMENT: "Problem 19. What can be said about the volume of the tetrahedron in the four-body problem? Can it be nonzero forever and stay bounded away from 0 and ∞? (Joseph Gerver)"
% Preparation and final self-audit: CONTRIBUTING.md prompts C and D; source/audit artifacts retained in the local physics_statements_2026_10_11 batch.
% !TeX program = lualatex
\documentclass[11pt]{article}
\usepackage[margin=25mm]{geometry}
\usepackage{fontspec}
\setmainfont{Latin Modern Roman}
\usepackage{amsmath,amssymb,mathtools}
\usepackage{unicode-math}
\setmathfont{Latin Modern Math}
\usepackage{xurl,hyperref}
\hypersetup{colorlinks=true,urlcolor=blue}
\setcounter{secnumdepth}{0}
\setlength{\parindent}{0pt}
\setlength{\parskip}{5pt}
\setlength{\emergencystretch}{3em}
\newcommand{\sref}[2]{\hyperlink{source:#1:#2}{[S#2]}}
\newcommand{\cataloguescope}[1]{\paragraph{Recorded target.}#1}
\newcommand{\R}{\mathbb R}
\newcommand{\C}{\mathbb C}
\newcommand{\N}{\mathbb N}
\newcommand{\Z}{\mathbb Z}
\newcommand{\Q}{\mathbb Q}
\newcommand{\F}{\mathbb F}
\newcommand{\eps}{\varepsilon}
\begin{document}
% UnsolvedMath ID 20002548; source catalogue code AIM-PHYSICS-0056
\setcounter{section}{20002547}
\section{A logarithmic-volume identity and zero-angular-momentum obstruction}\label{20002548}
{\small UnsolvedMath ID 20002548 · Mathematical Physics}
\subsection{Definitions and mathematical statement}

\paragraph{Shared notation.}$\mathbb N=\{1,2,\ldots\}$; the standard number fields have their usual meanings. All additional conventions are stated locally.


\paragraph{Definitions and notation.}In the Newtonian $n$-body problem, positions $q_i\in\mathbb R^d$ and masses $m_i>0$ obey $m_i\ddot q_i=\nabla_{q_i}U$, where $U(q)=\sum_{i<j}m_im_j/|q_i-q_j|$ and the gravitational constant is normalized to one. Here $T=\tfrac12\sum_i m_i|\dot q_i|^2$, $E=T-U$, and a collision has $q_i=q_j$. The center of mass can be fixed at zero; moment of inertia is $I=\sum_i m_i|q_i|^2$. For four spatial positions, tetrahedron volume is $V(t)=\tfrac16|\det(q_2-q_1,q_3-q_1,q_4-q_1)|$. Remaining bounded away from zero and infinity means that positive finite constants $v_-,v_+$ bound $V(t)$ throughout the stated time interval.

\paragraph{Catalogue wording retained for traceability.}
What restrictions govern the volume of the tetrahedron in the four-body problem? Can an orbit exist with $0<v_-\le V(t)\le v_+<\infty$ forever?
\paragraph{Scope and source qualification.}Forever must specify one or both directions of time and exclude breakdown by collision. The source does not add zero angular momentum or a chosen energy regime, so such conditions are not inserted. The catalogue title/status was paired with a machine-generated research attempt; that attempt is not peer-reviewed evidence of settlement. The statement covers the original source questions rather than only the title’s subcase.
\subsection{Short English statement}
Determine whether a four-body orbit can maintain a positive tetrahedron volume uniformly bounded above and below forever.
\subsection{Sources}
\begin{itemize}
\item[\hypertarget{source:20002548:1}{S1}] ULAM.AI and the UnsolvedMath Contributors, UnsolvedMath, record 20002548 (AIM-PHYSICS-0056). Catalogue curation under CC BY 4.0; definitions and English exposition added locally.
\url{https://huggingface.co/datasets/ulamai/UnsolvedMath}.
\item[\hypertarget{source:20002548:2}{S2}] American Institute of Mathematics, Variational Methods in Celestial Mechanics; notes by Kuo-Chang Chen, version 22 June 2003, source item 19.
\url{https://aimath.org/WWN/varcelest/varcelest.pdf}.
\end{itemize}
\label{end:20002548}
\end{document}
